%% kapa: exam=1 round=23 year=2012 period=2 form=A subject=appraisal_law subject_ko=부동산관계법규 seq=19 number=59
%% source: Q-Net 감정평가사 기출문제 (cst00309) artlSeq=5052580 file=raw/1008/5052580/2012년 제23회 감정평가사 제1차시험 기출문제(2교시 A형).pdf
%% conversion: kordoc 4.24.51 -> tools/gen1.py (auto); review=proofread-sonnet; flags=none
\begin{question}{59}
부동산 가격공시 및 감정평가에 관한 법령상 감정평가업자 또는 감정평가사에 관한 설명으로 옳지 않은 것은?
\begin{choices}
  \item 국토해양부장관은 감정평가사의 자격등록취소처분을 하고자 하는 경우 청문을 실시하여야 한다.
  \item 감정평가업자는 감정평가서의 원본을 교부일로부터 10년 이상 보존하여야 한다.
  \item 파산선고를 받고 복권된 자는 감정평가법인의 대표이사가 될 수 있다.
  \item 감정평가업자는 의뢰인에 대한 손해배상책임을 보장하기 위하여 보증보험이나 공제사업에 가입하고, 이를 국토해양부장관에게 통보하여야 한다.
  \item 부정한 방법으로 감정평가사의 자격을 얻었음을 이유로 그 자격이 취소된 후 3년이 경과되지 아니한 자는 감정평가사가 될 수 없다.
\end{choices}
\end{question}
\answer{1}
